\section{The exact swirl-rate and positive-carrier transfer}
\label{sec:swirl-carrier}

We retain the original axisymmetric coordinates of
Section~\ref{sec:exact-euler-coordinates}:
$y_1=z$, $y_2=r^2/2>0$, $v=(u^z,ru^r)=J\nabla\Psi$,
$\Gamma=ru^\phi$, and $\xi=L\Psi=\omega^\phi/r$.
The viscosity remains the fixed physical constant $\nu>0$.
The following change of dependent variable makes the meridional source
exactly a derivative of a square, while retaining the geometric terms
in its transport equation. It then gives a complete compact transfer
from a signed scalar layer into actual three-dimensional data and force.

\subsection{A common viscous operator for swirl rate and vorticity}

Define the signed angular velocity, or swirl rate, by
\[
 h=\frac{\Gamma}{2y_2}=\frac{u^\phi}{r},\qquad
 Q=\partial_1^2+2y_2\partial_2^2,\quad
 M=Q+4\partial_2,\quad D_v=\partial_t+v\cdot\nabla_y.
\]
This is an invertible change of variables on the retained domain
$y_2>0$: its inverse is $\Gamma=2y_2h$. No time or spatial coordinate
has been rescaled. Product differentiation gives
\begin{align}
 Q(2y_2h)&=2y_2h_{11}+2y_2(4h_2+2y_2h_{22})=2y_2Mh,
 \label{eq:sc-diffusion}\\
 D_v(2y_2h)&=2y_2D_vh+2v_2h,\qquad
 W\partial_1(\Gamma^2)=\partial_1(h^2),
 \quad W=(2y_2)^{-2}.
 \label{eq:sc-transport}
\end{align}
Thus the complete reduced system \eqref{eq:ec-reduced-NS} is
equivalent to
\begin{equation}\label{eq:sc-rate-system}
 \boxed{\begin{aligned}
 v&=J\nabla\Psi,\qquad \xi=L\Psi,\\
 D_vh&=\nu Mh-\frac{v_2}{y_2}h+\frac{f^\phi}{r},\\
 D_v\xi&=\nu M\xi+\partial_1(h^2)
                    +\curl_y(f^z,f^r/r).
 \end{aligned}}
\end{equation}
Both scalar equations now use the same operator $M$. The multiplication
term $-v_2h/y_2$ is part of the exact swirl equation. In original physical
coordinates it is $-2u^rh/r$, since $v_2=ru^r$ and $y_2=r^2/2$.

For scalar functions $a,b$ write
\[
 \langle\nabla a,\nabla b\rangle_Q=a_1b_1+2y_2a_2b_2,
 \qquad |\nabla a|_Q^2=\langle\nabla a,\nabla a\rangle_Q.
\]
The product rule, including the drift $4\partial_2$, is
\[
 M(ab)=aMb+bMa+2\langle\nabla a,\nabla b\rangle_Q.
\]
It follows either by differentiating twice in both coordinates, or by
noting that the drift obeys the first-order product rule and contributes
no additional quadratic term. Set $B=h^2$, a new square distinct from
the original $\Theta=\Gamma^2$ in \eqref{eq:ec-Theta}. Multiplying the
swirl equation by $2h$ gives, everywhere including $h=0$,
\begin{equation}\label{eq:sc-square-global}
 D_v B=\nu MB-2\nu|\nabla h|_Q^2
       -\frac{2v_2}{y_2}B+\frac{2h}{r}f^\phi.
\end{equation}
On a component where $h\ne0$, $B>0$ and
$|\nabla B|_Q^2=4B|\nabla h|_Q^2$. Hence there the same equation is
\begin{equation}\label{eq:sc-square-positive}
 D_v B=\nu MB-\frac{\nu}{2B}|\nabla B|_Q^2
       -\frac{2v_2}{y_2}B+\frac{2h}{r}f^\phi.
\end{equation}
This last formula is only used on $B>0$. The globally smooth equation
\eqref{eq:sc-square-global} is used at zeros, so no division by a
vanishing swirl has been inserted.

\subsection{A signed layer, positive carrier and explicit compact forces}

Retain a bounded convex ball $D\Subset\{y_2>0\}$ and the compact
curl inverse $\mathcal R$ constructed in
Section~\ref{sec:exact-euler-coordinates}. Fix $0<T_0<T_1$ and an
open time interval $I$ containing $[0,T_1]$. Take real smooth fields
$\Psi,\vartheta$ on $\R^2\times I$ whose spatial supports lie in fixed
compact subsets of $D$. Choose a smooth cutoff
$0\leq\chi\leq1$, compactly supported in $D$, equal to one on a
neighborhood of the spatial support of $\vartheta$ at every time.
To give an explicit choice, write $D=B_R(y_*)$ and let $K$ be the
union of the fixed compact support set with $\{y_*\}$. It is a
nonempty compact subset of $D$, including when $\vartheta=0$. Choose
$\max_{y\in K}|y-y_*|<R_0<R_1<R$ and use the spatial cutoff
in \eqref{eq:vs-cutoffs} with $|x|$ replaced by $|y-y_*|$.
Compactness of $K\Subset D$ guarantees these strict inequalities.

Retain an amplitude conversion factor $\mu>0$. Define
\[
 A_\vartheta=\sup_{\R^2\times[0,T_1]}|\vartheta|<\infty,
 \qquad c>\mu A_\vartheta,
 \qquad B_0=c+\mu\vartheta,
 \qquad H=\sqrt{B_0}>0.
\]
The finite supremum exists by smoothness on the fixed compact support
and compact time interval. Put $\delta=c-\mu A_\vartheta>0$.
The bound $B_0\geq\delta$ holds on $[0,T_1]$, and uniform continuity
on a slightly larger compact cylinder gives $B_0\geq\delta/2>0$
on some time neighborhood of $[0,T_1]$.
The positive root $H$ and all derivatives are consequently
smooth there. The coefficient $\mu$ is part of the scalar amplitude
dictionary: $\mu\vartheta$ has the units of the angular velocity
squared, as does $c$. No claim identifying unrelated scalar units is
implicit in this definition.

Let $\eta(t)$ be the explicit smooth time cutoff in
\eqref{eq:vs-cutoffs}, with these same $T_0,T_1$. Thus $\eta=1$
on $[0,T_0]$, $\eta=0$ on $[T_1,\infty)$, and every derivative
vanishes at $T_1$. Put
\begin{equation}\label{eq:sc-compact-fields}
 \widetilde\Psi=\eta\Psi,\quad
 v_0=J\nabla\Psi,\quad \xi_0=L\Psi,\quad
 v=\eta v_0,\quad\xi=\eta\xi_0,\quad
 h=\eta\chi H,\quad\Gamma=2y_2h.
\end{equation}
The tilde marks the actual target streamfunction; $\Psi$ is the
original input and is not silently replaced in a derivative.
Define two original-input residuals
\begin{align}
 E_\vartheta&=\partial_t\vartheta+v\cdot\nabla\vartheta-\nu M\vartheta,
 \label{eq:sc-input-scalar}\\
 E_0&=\partial_t\xi_0+v_0\cdot\nabla\xi_0-\nu M\xi_0
                   -\mu\partial_1\vartheta.
 \label{eq:sc-input-vorticity}
\end{align}
Their different velocities, $v$ and $v_0$, are deliberate and retained
in the following full formulas. Prescribe
\begin{align}
 R_h={}&\eta'\chi H
 +\eta\chi\left[
       \frac{\mu E_\vartheta}{2H}
       +\frac{\nu\mu^2}{4H^3}|\nabla\vartheta|_Q^2
       +\frac{v_2}{y_2}H\right]\notag\\
 &+\eta H(v\cdot\nabla\chi-\nu M\chi)
       -\frac{\eta\nu\mu}{H}
                    \langle\nabla\chi,\nabla\vartheta\rangle_Q,
 \label{eq:sc-swirl-residual}\\
 E_\xi={}&\eta E_0+\eta'\xi_0
 +\eta(\eta-1)
           (v_0\cdot\nabla\xi_0-\mu\partial_1\vartheta)
       -\eta^2c\partial_1(\chi^2).
 \label{eq:sc-vorticity-residual}
\end{align}
In the actual fixed-support construction the last inner product in
\eqref{eq:sc-swirl-residual} is zero, because $\chi$ is identically
one near the support of $\vartheta$. It is retained to expose the
complete product rule and the additional term for a nonflat envelope.

\begin{proposition}[Complete finite-interval transfer]
\label{prop:sc-transfer}
For the explicitly defined fields above, put
\begin{equation}\label{eq:sc-forces}
 b=\mathcal R E_\xi,\qquad
 f^z=b_1,\quad f^r=rb_2,\quad f^\phi=rR_h.
\end{equation}
The physical velocity with streamfunction $\widetilde\Psi$ and
circulation $\Gamma$ solves the original three-dimensional viscous
equation for a smooth pressure. Its data and force satisfy
\eqref{eq:schwartz-data}--\eqref{eq:schwartz-force}, and its energy
has a uniform bound for all $t\geq0$.
On $0\leq t\leq T_0$ and the plateau $\chi=1$, the exact identities are
\[
 h^2=c+\mu\vartheta,\qquad
 \partial_1(h^2)=\mu\partial_1\vartheta,\qquad
 E_\xi=E_0,
\]
and
\begin{equation}\label{eq:sc-core-force}
 \frac{f^\phi}{r}
 =\frac{\mu E_\vartheta}{2\sqrt{c+\mu\vartheta}}
  +\frac{\nu\mu^2|\nabla\vartheta|_Q^2}
                    {4(c+\mu\vartheta)^{3/2}}
  +\frac{v_2}{y_2}\sqrt{c+\mu\vartheta}.
\end{equation}
\end{proposition}
\begin{proof}
The chain rule with $H=\sqrt{c+\mu\vartheta}$ gives
\[
 D_vH-\nu MH
 =\frac{\mu E_\vartheta}{2H}
   +\frac{\nu\mu^2}{4H^3}|\nabla\vartheta|_Q^2,
 \qquad \nabla H=\frac\mu{2H}\nabla\vartheta.
\]
For verification, the first two derivatives of the function
$s\mapsto\sqrt{c+\mu s}$ are $\mu/(2\sqrt{c+\mu s})$
and $-\mu^2/(4(c+\mu s)^{3/2})$. Inserting them into the
second-order part of $M$ supplies its displayed positive residual
term; the drift obeys the ordinary first derivative rule.
Apply the product rule to $h=\eta\chi H$, remembering that
$D_v\eta=\eta'$, $D_v\chi=v\cdot\nabla\chi$, and $\eta$
is constant in space. The result is exactly
\[
 D_vh-\nu Mh+(v_2/y_2)h=R_h
\]
with every summand in \eqref{eq:sc-swirl-residual}.

The plateau condition gives the global identity
$\chi^2\vartheta=\vartheta$, and therefore
\[
 h^2=\eta^2(c\chi^2+\mu\vartheta),\qquad
 \partial_1(h^2)=\eta^2
                    [c\partial_1(\chi^2)+\mu\partial_1\vartheta].
\]
Next $D_v\xi=\eta'\xi_0+\eta\partial_t\xi_0
+\eta^2v_0\cdot\nabla\xi_0$ and $M\xi=\eta M\xi_0$.
Subtracting the preceding source and $\nu M\xi$ gives
\eqref{eq:sc-vorticity-residual} exactly, so
$D_v\xi-\nu M\xi-\partial_1(h^2)=E_\xi$.

This residual has zero spatial integral. Indeed
$\int\xi=\int L\widetilde\Psi=0$, since integration of each
second derivative in its own coordinate gives zero and the coefficient
of $\partial_1^2$ is independent of $y_1$.
Divergence-free $v$ gives $\int v\cdot\nabla\xi=0$;
$\int\partial_1(h^2)=0$; and
$\int M\xi=0$ by the explicit integration by parts already proved
after \eqref{eq:ec-residuals}. Thus the compact inverse applies and
$\curl_y b=E_\xi$. Equations \eqref{eq:sc-forces} satisfy both
equations in \eqref{eq:sc-rate-system}. Equivalently
$F_\nu^\Gamma=2y_2R_h$ and the force is precisely the force in
Proposition~\ref{prop:ec-reconstruction}. That proposition supplies
the full velocity and the explicitly integrated smooth pressure, so
no pressure compatibility or missing vector component is assumed here.

All input supports and the cutoff support are fixed compact subsets
of $D$. Division by $H$ is uniformly harmless because of its positive
lower bound, and $r$ has a fixed positive lower bound there. The
compact inverse has the smoothness and support proved in that
proposition. Every displayed field is therefore smooth on a time
neighborhood of $[0,T_1]$ within its support. At $T_1$, $\eta$ and
all its derivatives vanish; the formulas for $E_\xi$ and $R_h$
show the same flatness for the force. Extending the velocity, force
and the reconstructed pressure by zero after $T_1$ is smooth.
The velocity and force have spatial support in the fixed solid torus
corresponding to $D$. Hence all initial-data derivatives are rapidly decaying in
space. Every force derivative is bounded on its compact spacetime
support, which proves \eqref{eq:schwartz-force} for every stated
multi-index, time derivative and polynomial weight.

For an explicit energy bound let
$y_- =\inf_D y_2>0$, $y_+=\sup_D y_2<\infty$, and let
$S_j=\sup_{D\times[0,T_1]}|\partial_j\Psi|$.
The energy identity proved in the coordinate section becomes exactly
\[
 \int_{\R^3}|u|^2\,dx
 =2\pi\eta^2\int_D
 \left[\Psi_2^2+\frac{\Psi_1^2}{2y_2}
                    +2y_2\chi^2(c+\mu\vartheta)\right]dy.
\]
It is at most
$2\pi |D|[S_2^2+S_1^2/(2y_-)+2y_+(c+\mu A_\vartheta)]$
on $[0,T_1]$, and is zero after $T_1$. This finite number plus one
supplies \eqref{eq:energy}. The plateau formulas follow by setting
$\eta=\chi=1$ and all their derivatives to zero in the exact residuals.
\end{proof}

The carrier changes the meridional source only by the explicit fixed
edge term $c\partial_1(\chi^2)$ during the uncut time interval.
Its swirl force retains both the geometric term and the positive
viscosity gradient-square term in \eqref{eq:sc-core-force}.
These formulas are an exact reconstruction map with all forces
computed. They construct global smooth solutions for fixed smooth
inputs and finite cutoffs; they do not establish a singular limit.
